\chapter{Metodología}
\label{cap:metodologia}
En este capítulo se presenta el marco metodológico empleado para analizar la estructura de dependencia entre activos financieros mediante técnicas de \gls{HRP} y su extensión cuántica. Se estudian las relaciones latentes del mercado utilizando métricas alternativas a la correlación tradicional, permitiendo una representación más rica y robusta de las interdependencias entre activos.

\section{Representación cuántica de activos financieros}
\label{sec:representacion_cuantica_activos}

\subsection{Representación multivariante de los activos}
\label{subsec:representacion_multivariante}

Cada activo financiero se representa mediante un conjunto de características económicas que describen distintos aspectos de su comportamiento dinámico. Para cada activo $i$ se construye una matriz de observaciones:

\begin{equation}[eq:matriz_observaciones]{Matriz de observaciones de un activo individual}
X_i \in \mathbb{R}^{T \times P},
\end{equation}

donde $T$ denota el número de observaciones temporales y $P$ el número de características consideradas. En este TFG se emplean seis características por activo, calculadas a partir de ventanas móviles de retornos: (1) retorno diario, (2) media móvil de retornos, (3) desviación estándar móvil, (4) retorno acumulado, (5) asimetría (\textit{skewness}) y (6) curtosis. Esta combinación de momentos estadísticos de la distribución de retornos permite capturar tanto el comportamiento central del activo como sus características de cola y dispersión, proporcionando una representación multifacética del perfil de riesgo.

El uso de una representación multivariante permite superar las limitaciones del enfoque clásico basado exclusivamente en la correlación de retornos, la cual es conocida por su elevada inestabilidad temporal y sensibilidad al ruido.

\subsection{Normalización de características}
\label{subsec:normalizacion_caracteristicas}

La adecuación de las características financieras al dominio cuántico se realiza mediante una normalización angular integrada en el mapa de características (en inglés, \textit{feature map}). Cada una de las seis variables $x_i$ se transforma, mediante un mapeo afín de tipo min-max respecto a sus extremos históricos en la ventana temporal activa, en un ángulo $\theta_i \in [0, \alpha\pi]$, que actúa directamente como parámetro de las puertas de rotación cuánticas $R_y$ y $R_x$ del circuito (véase Apéndice~\ref{ap:subsubsec:puertasunqubit} para la definición de las puertas de rotación y Apéndice~\ref{ap:estimacion_ansatz} para los detalles del ansatz).
Este procedimiento permite que el codificación cuántica se adapte dinámicamente a los cambios de régimen de mercado sin requerir estandarizaciones clásicas adicionales. Al mapear los datos al rango operativo de las puertas de rotación, se asegura que la información quede codificada correctamente en el espacio de Hilbert para su posterior análisis.


\subsection{Representación cuántica mediante matrices de densidad}
\label{subsec:representacion_cuantica}
Para modelar relaciones no lineales entre activos, cada observación multivariante se codifica en un estado cuántico mediante un mapa de características parametrizado. A partir de dicho mapeo, cada observación queda representada como un estado cuántico puro:

\begin{equation}[eq:estado_cuantico_observacion]{Estado cuántico de una observación de un activo individual}
\lvert\psi(\mathbf{x}_i^t)\rangle.
\end{equation}

Posteriormente, para cada activo se construye una matriz de densidad promedio definida como:

\begin{equation}[eq:matriz_densidad_promedio]{Matriz de densidad promedio de un activo}
\rho_i = \frac{1}{T} \sum_{t=1}^{T}
\lvert\psi(\mathbf{x}_i^t)\rangle\langle\psi(\mathbf{x}_i^t)\rvert.
\end{equation}

Esta representación mediante matrices de densidad permite que el modelo sea sensible no solo al valor medio de las características, sino a la distribución completa y estructura de la nube de puntos en el espacio de Hilbert. Al tratar al activo como un estado mixto de sus observaciones históricas, se proporciona una representación más informativa que los estadísticos clásicos de segundo orden, capturando dependencias que la correlación lineal ignora.

El mapa de características puede incluir capas de entrelazamiento mediante puertas de dos qubits (en este caso iSWAP), lo que permite capturar dependencias no lineales en el espacio de Hilbert.

\subsection{Métrica de distancia cuántica}
\label{subsec:metrica_distancia_cuantica}
La similitud entre activos se cuantifica mediante la distancia de Frobenius entre sus matrices de densidad promedio. Para dos activos $i$ y $j$, dicha distancia se define como:

\begin{equation}[eq:distancia_frobenius_cuantica]{Distancia de Frobenius entre matrices de densidad}
d_F(\rho_i, \rho_j) = \frac{1}{2}
\sqrt{\mathrm{Tr}\left[(\rho_i - \rho_j)^2\right]}.
\end{equation}

Esta métrica proporciona una medida geométrica natural en el espacio de operadores cuánticos y permite comparar activos teniendo en cuenta la totalidad de su comportamiento multivariante.

\subsection{Construcción de la matriz de distancias y ordenamiento jerárquico}
\label{subsec:construccion_ordenamiento}

Las distancias cuánticas entre todos los pares de activos se organizan en una matriz de distancias simétrica $D \in \mathbb{R}^{N \times N}$. Dicha matriz se emplea como entrada para un procedimiento de \textit{clustering} jerárquico, que genera un dendrograma jerárquico a partir de una métrica de distancia entre \textit{clusters}.

En la implementación experimental comparada en este TFG, el \gls{HRP} clásico y el \gls{KHRP} emplean enlace simple para la construcción jerárquica, mientras que la ordenación cuántica evaluada con distancias de Frobenius utiliza Ward. El \textbf{enlace simple} (en inglés \textit{single linkage}) mide la distancia entre dos grupos como la distancia entre los dos activos más cercanos de cada grupo, lo que puede dar lugar a estructuras alargadas en el dendrograma. El \textbf{criterio de Ward} en cambio fusiona en cada paso los dos grupos cuya unión provoca el menor aumento de varianza interna, tendiendo a formar grupos más compactos y equilibrados. Esta diferencia afecta únicamente a la fase de ordenamiento; el esquema de asignación de pesos mediante bisección recursiva se mantiene común.

Del dendrograma se extrae un ordenamiento de los activos, de forma que activos con menor distancia tienden a situarse próximos en la representación final. El reordenamiento de la matriz de distancias mediante este procedimiento no altera los valores originales, pero facilita su reorganización consistente con la estructura jerárquica.

\section{Métricas de evaluación del rendimiento}
\label{sec:metricas_evaluacion}
Para evaluar y comparar el rendimiento de las estrategias de inversión resultantes de los enfoques clásico y cuántico, se emplearán métricas estandarizadas que consideran tanto la rentabilidad obtenida como el riesgo asumido.

\subsection{Sharpe Ratio}
\label{subsec:sharpe_ratio}
El \textbf{Sharpe Ratio} es la métrica de referencia en entidades financieras para medir el rendimiento de una cartera ajustado por riesgo. Representa el exceso de retorno por unidad de desviación estándar (volatilidad).

\begin{equation}[eq:sharpe_ratio]{Sharpe Ratio}
    SR = \frac{E[R_p - R_f]}{\sigma_p}
\end{equation}

Donde $R_p$ es el retorno de la cartera, $R_f$ es la tasa libre de riesgo y $\sigma_p$ es la volatilidad de los retornos de la cartera.

\begin{itemize}
    \item \textbf{Interpretación:} Un Sharpe Ratio más alto indica una gestión más eficiente, sugiriendo que la rentabilidad no es fruto de una exposición excesiva al riesgo.
    \item \textbf{Limitaciones:} Su principal debilidad es que asume que los retornos siguen una distribución normal. En mercados financieros, donde existen eventos de cola ancha (\textit{fat tails}) y asimetría, el Sharpe Ratio puede subestimar el riesgo real.
\end{itemize}

\subsection{PSR}
\label{subsec:probabilistic_sharpe_ratio}
Dado que el Sharpe Ratio es una estimación estadística sujeta a ruido, se utilizará el \textbf{\gls{PSR}} para evaluar la significación de los resultados. \gls{PSR} corrige los sesgos introducidos por la asimetría (en inglés \textit{skewness}) y la curtosis de la distribución de retornos, calculando la probabilidad de que el Sharpe Ratio observado sea superior a un valor de referencia (benchmark). Esta métrica es fundamental en este TFG para validar si la mejora del \gls{QHRP} es estadísticamente robusta o producto del azar.

\subsection{Otras métricas de riesgo}
\label{subsec:otras_metricas_riesgo}
Complementariamente, se incorpora la métrica \textbf{\gls{MinTRL}} para las variantes \gls{HRP}, como indicador de robustez estadística de la señal. En esta implementación, el análisis cuantitativo principal de rendimiento se centra en Sharpe y \gls{PSR}; la estabilidad se interpreta a partir de la dispersión interventana de dichas métricas.